\BeleskeID{04-s036}
% element: 04-s036:e01 — Negativna vrednost i komplementiranje sa minusom
% element: 04-s036:e02 — Početak izvođenja i rezultat preko bita najveće težine
% element: 04-s036:e03 — Uslov bn−1=1 i negativna težina vodećeg bita
% element: 04-s036:e04 — Naglašen zapis (−bn−1)bn−2…b0
% element: 04-s036:e05 — Dva algebarska međukoraka izdvajanja faktora 2^(n−1)
% element: 04-s036:d01

Niz $b_{n-2}b_{n-3}\ldots b_1b_0$ u računu označava celobrojnu vrednost donjih $n-1$ bitova, a ne njihov proizvod. Puno izvođenje glasi:
\[\begin{aligned}
&-\left(2^n-b_{n-1}b_{n-2}b_{n-3}\ldots b_1b_0\right)\\
&=-\left(2^n-b_{n-1}2^{n-1}-b_{n-2}b_{n-3}\ldots b_1b_0\right)\\
&=-\left((2-b_{n-1})2^{n-1}-b_{n-2}b_{n-3}\ldots b_1b_0\right)\\
&=-(2-b_{n-1})2^{n-1}+b_{n-2}b_{n-3}\ldots b_1b_0\\
&=(b_{n-1}-2)2^{n-1}+b_{n-2}b_{n-3}\ldots b_1b_0.
\end{aligned}\]
Za negativni kod $b_{n-1}=1$, pa je koeficijent vodeće težine $1-2=-1=-b_{n-1}$. Donji biti zadržavaju pozitivne težine.

Za kod $1101$ u četiri bita čitamo $-8+4+0+1=-3$. Skraćeni zapis sa $(-b_{n-1})$ naglašava negativnu težinu vodećeg bita; on ne znači da fizička ćelija registra čuva „negativnu jedinicu“. Ćelija i dalje čuva običan bit $1$.
